\documentclass[12pt]{article}
\usepackage[a4paper,margin=1in]{geometry}\usepackage[T1]{fontenc}
\usepackage{lmodern,microtype}\usepackage{amsmath,amssymb,amsthm,mathtools}
\usepackage{booktabs,array,enumitem,xcolor}\usepackage[numbers,sort&compress]{natbib}
\usepackage[colorlinks=true,linkcolor=blue!55!black,citecolor=blue!55!black,urlcolor=blue!55!black]{hyperref}
\linespread{1.07}
\newtheorem{theorem}{Theorem}[section]\newtheorem{proposition}[theorem]{Proposition}
\theoremstyle{remark}\newtheorem{remark}[theorem]{Remark}
\newcommand{\TV}{\operatorname{TV}}

\title{A Finite Monotone Axial Clock for the Quartet Shape Flow\\
\large Sixteen Scales, Two Channels, and a Narrow Asymmetry Corridor}
\author{Liang Wang}\date{July 2026}
\begin{document}\maketitle
\begin{abstract}
RH-213 reduces each centered physical quartet to exact shape coordinates
$(u,\eta)\in[0,1]\times[-1,1]$.  We now audit their scale ordering on sixteen
small-noise levels from $0.04$ to $0.00125$.  On both physical channels, $u$
increases strictly through all fifteen coarse-to-fine transitions.  The net
change is $0.69809$ on the left and $0.69673$ on the right; the smallest
observed increments are $0.01529$ and $0.01513$.

After the initial transient, defined in advance here by $\sigma\le0.02$,
$\eta$ remains in narrow corridors: $[-0.10444,-0.06886]$ on the left and
$[-0.10748,-0.07031]$ on the right.  Their widths are $0.03557$ and
$0.03717$.  Mature left/right discrepancies are at most $0.003006$ in $u$
and $0.003043$ in $\eta$.

Any strictly ordered finite data define an invertible piecewise-linear clock
on their sampled log-scale interval; we state this elementary theorem to
separate a legitimate finite clock from an unsupported all-level law.  The
observed trajectory is therefore well described as one dominant axial motion
plus a bounded transverse corridor, but no monotonicity theorem below the
finest level, autonomous flow, or small-noise limit is claimed.
\end{abstract}

\section{Exact coordinates and finite question}

For each physical channel $s\in\{L,R\}$ and scale $\sigma$, RH-213 writes
the centered-RMS quartet as
\begin{equation}\label{eq:roots}
 \sqrt{u_{\sigma,s}}\pm i\sqrt{(1-u_{\sigma,s})(1+\eta_{\sigma,s})},
 \quad
 -\sqrt{u_{\sigma,s}}\pm i\sqrt{(1-u_{\sigma,s})(1-\eta_{\sigma,s})}
\end{equation}
\citep{WangRH213}.  The coefficient vector is an exact function of these two
coordinates.  RH-212 supplies a frozen sixteen-level root atlas
\citep{WangRH212}.

We ask three finite questions:
\begin{enumerate}[leftmargin=2em]
\item Is $u$ ordered consistently as $\sigma$ decreases?
\item Does $\eta$ enter a narrower post-transient corridor?
\item Do both physical channels track the same coordinate path within small
finite discrepancy?
\end{enumerate}
None of these questions presupposes a limiting value.

\section{Scale order and clock convention}

Write
\begin{equation}\label{eq:time}
 t=\log(1/\sigma).
\end{equation}
The sixteen anchors are ordered by increasing $t$, equivalently decreasing
$\sigma$.  Let $t_0<\cdots<t_{15}$ and let $u_k=u_{\sigma_k,s}$ on one
channel.

\begin{theorem}[Finite monotone-clock interpolation]\label{thm:clock}
If
\begin{equation}\label{eq:strict}
 u_0<u_1<\cdots<u_n,
\end{equation}
then there is a unique continuous function $U:[t_0,t_n]\to[u_0,u_n]$ that
is affine on each interval $[t_k,t_{k+1}]$ and satisfies $U(t_k)=u_k$.
It is strictly increasing and hence a homeomorphism onto its image.
\end{theorem}
\begin{proof}
Linear interpolation is unique on every subinterval.  Its slope there is
$(u_{k+1}-u_k)/(t_{k+1}-t_k)>0$, so the pieces join to a strictly increasing
continuous bijection.  A continuous strictly monotone map on a compact
interval has continuous inverse.
\end{proof}

The theorem is intentionally finite.  It neither extends $U$ beyond the
sampled interval nor says that the piecewise-linear interpolation is the
physical law.

\section{Sixteen-level axial verdict}

For each transition define
\begin{equation}
 \Delta u_{k,s}=u_{k+1,s}-u_{k,s}.
\end{equation}
All thirty channel-transition increments are positive.

\begin{center}
\begin{tabular}{lrr}
\toprule
axial diagnostic & left & right\\
\midrule
transition count & 15 & 15\\
positive increments & 15 & 15\\
minimum increment & $0.015295$ & $0.015130$\\
maximum increment & $0.104022$ & $0.106002$\\
net change & $0.698086$ & $0.696733$\\
total variation & $0.698086$ & $0.696733$\\
rank correlation with $\log(1/\sigma)$ & $1$ & $1$\\
\bottomrule
\end{tabular}
\end{center}

Equality of total variation and net change is simply another expression of
finite monotonicity.  The factor-of-seven spread between minimum and maximum
increments warns against treating the anchor index itself as constant-speed
time.

\section{Representative dyadic subsequence}

The following equal-log-step subsequence displays the dominant motion:
\begin{center}
\small
\begin{tabular}{rrrrr}
\toprule
$\sigma$ & $u_L$ & $u_R$ & $\eta_L$ & $\eta_R$\\
\midrule
$0.04$ & $0.02027$ & $0.02077$ & $-0.57806$ & $-0.59358$\\
$0.02$ & $0.24887$ & $0.24794$ & $-0.06886$ & $-0.07031$\\
$0.01$ & $0.40529$ & $0.40229$ & $-0.10443$ & $-0.10748$\\
$0.005$ & $0.54098$ & $0.54034$ & $-0.08934$ & $-0.09080$\\
$0.0025$ & $0.65874$ & $0.65801$ & $-0.09050$ & $-0.09181$\\
$0.00125$ & $0.71835$ & $0.71750$ & $-0.09303$ & $-0.09444$\\
\bottomrule
\end{tabular}
\end{center}
The first transition contains a large transverse relaxation.  Thereafter the
axial coordinate advances while $\eta$ fluctuates around roughly $-0.09$.

\section{The mature asymmetry corridor}

For a finite ordered sequence $x_k$, define
\begin{equation}
 \operatorname{width}(x)=\max_kx_k-\min_kx_k,
 \qquad
 \TV(x)=\sum_k|x_{k+1}-x_k|.
\end{equation}
On the levels with $\sigma\le0.02$:
\begin{center}
\begin{tabular}{lrr}
\toprule
transverse diagnostic & left & right\\
\midrule
minimum $\eta$ & $-0.104434$ & $-0.107477$\\
maximum $\eta$ & $-0.068864$ & $-0.070311$\\
width & $0.035570$ & $0.037166$\\
mean & $-0.089618$ & $-0.091187$\\
RMS about mean & $0.009378$ & $0.009636$\\
total variation & $0.070188$ & $0.074222$\\
\bottomrule
\end{tabular}
\end{center}

The total variation is about twice the width, so $\eta$ is not monotone.  A
``transverse fixed point'' would overstate the evidence.  The justified term
is a finite corridor.

\section{Dual-channel synchronization}

At fixed $\sigma$, define
\begin{equation}
 \delta_u(\sigma)=|u_{\sigma,L}-u_{\sigma,R}|,
 \qquad
 \delta_\eta(\sigma)=|\eta_{\sigma,L}-\eta_{\sigma,R}|.
\end{equation}
Across all sixteen scales,
\begin{equation}
 \max\delta_u=0.004229,
 \qquad \max\delta_\eta=0.015523,
\end{equation}
where the transverse maximum occurs in the initial transient.  On the mature
subset,
\begin{equation}\label{eq:maturesync}
 \max\delta_u=0.003006,
 \qquad \max\delta_\eta=0.003043.
\end{equation}

These discrepancies are small compared with the net axial displacement, but
they are not interval-certified errors.  They quantify agreement of two
finite discretizations/channels under the same root-selection rule.

\section{A near-one-dimensional description}

The phrase ``near one-dimensional'' has a precise finite meaning here:
\begin{enumerate}[leftmargin=2em]
\item one coordinate is strictly ordered through every observed transition;
\item the other remains in a post-transient interval of width below $0.038$;
\item both channels agree to roughly $3\times10^{-3}$ in each mature
coordinate.
\end{enumerate}
It does not mean that $\eta$ is a function of $u$ at all levels, that the
trajectory lies on an invariant curve, or that a one-dimensional generator
has been found.

\section{Why finite monotonicity matters}

The raw quartic coefficient flow looked irregular because location, radius,
axial separation, and transverse asymmetry were mixed.  The exact shape
coordinates show that one component has a stable order.  This gives a useful
clock for:
\begin{enumerate}[leftmargin=2em]
\item predeclared out-of-sample extrapolation;
\item boundary-distance estimates;
\item sensitivity decompositions;
\item equal-log-step recurrence audits.
\end{enumerate}
These uses remain finite unless accompanied by uniform estimates.

\section{Claim boundary and next test}

The data do not prove:
\begin{itemize}[leftmargin=2em]
\item $u_\sigma$ is monotone for every sufficiently small $\sigma$;
\item $u_\sigma$ has a limit, much less limit one;
\item $\eta_\sigma$ converges;
\item a scale-independent map advances the shape;
\item a growing determinant exists.
\end{itemize}

RH-215 therefore freezes seven fine levels as training data and withholds two
finer levels for prediction.  Gate A stays open, and no statement about
Gates B--E or zeta zeros is made.

\bibliographystyle{plainnat}\bibliography{references}\end{document}
